%%%%%%%%%%%%%%%%%
% Dies ist ein Beispiel-Lebenslauf, erstellt mit altacv.cls (v1.6.4, 13. November 2021) von
% LianTze Lim (liantze@gmail.com), basierend auf dem
% BusinessInsider-Lebenslauf unter [http://www.businessinsider.my/a-sample-resume-for-marissa-mayer-2016-7/?r=US&IR=T](http://www.businessinsider.my/a-sample-resume-for-marissa-mayer-2016-7/?r=US&IR=T)
%
%% Dieser Lebenslauf kann unter den Bedingungen der LaTeX Project Public License, Version 1.3,
%% verteilt oder modifiziert werden, oder (nach Ihrer Wahl) einer späteren Version.
%% Die neueste Version dieser Lizenz finden Sie unter
%% [http://www.latex-project.org/lppl.txt](http://www.latex-project.org/lppl.txt)
%% und Version 1.3 oder höher ist in allen LaTeX-Distributionen
%% ab 01.12.2003 oder später enthalten.
%%%%%%%%%%%%%%%%

%% Verwenden Sie die Option "normalphoto", wenn Sie ein reguläres Foto anstelle eines kreisförmigen möchten
\documentclass[9pt,a4paper,withhyper]{altacv}

%% AltaCV verwendet das Paket fontawesome5.
%% Siehe [http://texdoc.net/pkg/fontawesome5](http://texdoc.net/pkg/fontawesome5) für eine vollständige Liste der Symbole.

% Passen Sie das Seitenlayout bei Bedarf an
\geometry{left=1.25cm,right=1.25cm,top=1.45cm,bottom=1.5cm,columnsep=1cm}

% Das Paket paracol ermöglicht das Festlegen von Textspalten parallel
\usepackage{paracol}
\usepackage{xcolor}
\pagecolor{tan!5!white}

% Ändern Sie die Schriftart je nachdem, ob
% Sie pdflatex oder xelatex/lualatex verwende
\iftutex
  % Bei Verwendung von xelatex oder lualatex:
  \setmainfont{Lato}
\else
  % Bei Verwendung von pdflatex:
  \usepackage[default]{lato}
\fi

% Ändern Sie die Farben bei Bedarf
\definecolor{umber}{HTML}{352315}
\definecolor{tawny}{HTML}{80471C}
\definecolor{syrup}{HTML}{481F01}
\colorlet{heading}{syrup}
\colorlet{headingrule}{syrup}
\colorlet{accent}{tawny}
\colorlet{emphasis}{umber}
\colorlet{body}{umber}

% Ändern Sie bei Bedarf einige Schriften
% \\renewcommand{\\namefont}{\\Huge\\rmfamily\\bfseries}
% \\renewcommand{\\personalinfofont}{\\footnotesize}
% \\renewcommand{\\cvsectionfont}{\\LARGE\\rmfamily\\bfseries}
% \\renewcommand{\\cvsubsectionfont}{\\large\\bfseries}

% Ändern Sie die Symbole für itemize und Bewertungsmarker
% bei \\cvskill bei Bedarf
\renewcommand{\itemmarker}{{\small\textbullet}}
\renewcommand{\ratingmarker}{\faCircle}

\begin{document}
\name{Anirudh Singh}
\photoR{4cm}{static/profile.png}
\personalinfo {
  \extrainfo{\textcolor{accent}{\faGlobe} \href{https://praani.dev}{praani.dev}}
  \email{anirudh@mailbox.org}
  \github{ani-4nirudh}
  \linkedin{anirudhsingh8296}
}

\makecvheader

\columnratio{0.55}

\begin{paracol}{2}

\cvsection{Professional Experience}
  \cvevent{Software Support Engineer}{NEURA Robotics}{November 2025 -- April 2026}{Metzingen, Deutschland}
    \begin{itemize}
      \item Supporting customer service with complex technical issues
      \item Collaborating with development teams on testing, reproduction, and debugging of software features
      \item Contributing to the testing and development of new features across the full software stack
      \item Developing internal tools: Robot Deployment Virtual Machine Tool, logger tool for robot motion benchmarking, automation tools
    \end{itemize}

  \divider

  \cvevent{Support \& Applications Engineer}{Cambrian Robotics GmbH}{March 2024 -- August 2025}{Augsburg, Germany}
    \begin{itemize}
        \item Programming industrial and collaborative robots for feasibility studies  
        \item Developing and optimizing APIs for robotics software  
        \item Performing development testing, debugging, and validation  
        \item Technical customer support, including robot programming and automation consulting  
        \item Project management from concept phase to implementation  
        \item Developing and integrating mechatronic systems, including sensor integration with robot controllers (GPIO, PCB design)  
        \item Creating CAD models and 3D-printed components for custom hardware solutions  
    \end{itemize}

  \divider

  \cvevent{Master Thesis}{RWTH Aachen LLT - Chair of Laser Technology}{March 2023 -- December 2023}{Aachen, Deutschland}
    \begin{itemize}
        \item Goal: Identifying inaccuracies in robot TCP (Tool Center Point) position and velocity measurements using an onboard-mounted sensor
        \item Developing a laser-camera sensor concept for position estimation using image processing on laser-speckle images
        \item Programming classical image processing techniques in C++
        \item Identifying and testing the effectiveness of quantitative image-quality metrics for laser-speckle images
        \item Grade: 1.3
    \end{itemize}

\cvsection{Skills}
  \cvtag{Linux}
  \cvtag{POSIX}
  \cvtag{Embedded C}
  \cvtag{Real-time}
  \cvtag{Flask API}
  \cvtag{Git}
  \cvtag{Python}
  \cvtag{C++}
  \cvtag{Hardware Integration}
  \cvtag{Image Processing}
  \cvtag{Bash Scripting}
  \cvtag{OpenCV}
  \cvtag{ROS 2}
  \cvtag{Gazebo}

% \divider

% \cvevent{Working Student}{RWTH Aachen IP - Chair of Individualized Production}{May 2022 -- February 2023}{Aachen, Deutschland}
% \begin{itemize}
% \item Training a YOLO neural network for detecting personal protective equipment (PPE) using a custom dataset through transfer learning
% \item Building foundational knowledge of Linux networking and setting up ROS networks for communication over a 5G router (master and slave)
% \item Using RADAR, LiDAR, and 3D cameras together with ROS 1 (Noetic) for SLAM (Simultaneous Localization and Mapping)
% \item Implementing visual SLAM using ORB-SLAM3 and RTAB-Map packages
% \item Configuring a standalone system using a Jetson Xavier NX for remote control of an Innok robot over a 5G network
% \end{itemize}
%
% \divider
%
% \cvevent{Tutor - ROS Summer School 2022}{Fachhochschule Aachen}{15. Aug 2022 -- 29. Aug 2022}{Aachen, Deutschland}
%
% \begin{itemize}
% \item Week 1: Taught the basics of ROS2 (Gazebo, AprilTag detection)
% \item Week 2: Taught the basics of ROS2 navigation, localization, and mapping
% \end{itemize}
%
% \cvsection{Projects}
% \cvevent{Computer Vision and Deep Learning with Nvidia Jetson Nano}{Fachhochschule Aachen}{}{}
% \begin{itemize}
% \item Object detection and tracking with a mounted camera on servomotors
% \item Face recognition and face detection with the jetson-inference NVIDIA library and OpenCV
% \item Training a deep neural network with transfer learning
% \item Object tracking with contour detection and HSV color space
% \item Basic image processing techniques (thresholding, masking)
% \end{itemize}
%
% \divider
%
% \cvevent{Robotics Art Gallery}{Fachhochschule Aachen}{}{}
% \begin{itemize}
% \item Inspection of an "art gallery" by a robot in a Gazebo-ROS simulation
% \item Object detection with OpenVINO and YOLOX NN ROS package
% \item Autonomous navigation of the TurtleBot in the simulation
% \item Implementation of AprilTag detection
% \end{itemize}
%
% \divider

% \cvevent{Softwarearchitektur für autonomen elektrischen Kart}{FH Aachen}{}{}
% \begin{itemize}
% \item Verwendung von LIDAR mit ROS 2 (Foxy) für Mapping
% \item Verwendung einer 3D-Kamera für Objekt- und Feature-Detektion mit ROS 2 (Foxy)
% \item Sensorfusion von LIDAR und 3D-Kamera für SLAM
% \end{itemize}
%
% \divider
%
% \cvevent{Linienfolger-Roboter}{FH Aachen}{}{}
% \begin{itemize}
% \item Verwendung von Servos, Ultraschall zur Kollisionserkennung
% \item Infrarotsensoren zur Linienerkennung
% \item Entwurf eines Webservers zur Fernsteuerung des Roboters mit ESP32 und Arduino Uno
% \end{itemize}
%
\switchcolumn

\cvsection{Education}
  \cvevent{M.Sc. in Mechatronics}{Fachhochschule Aachen}{October 2020 -- December 2023}{}
  \begin{itemize}
    \item \href{https://drive.google.com/file/d/1J5xoOq-rh2UYf5aAfnolYBZ7_LzZLqmU/view}{Master’s thesis: Development of a Laser-Speckle-Based Velocity Sensor for Robot Tool Center Point Motion}
    \item Final grade: 1.7
  \end{itemize}

  \divider

  \cvevent{B.Sc. in Automobile Engineering}{Vellore Institute of Technology, Vellore}{June 2014 -- May 2018}{}
  \begin{itemize}
      \item Bachelor’s thesis: Design and Development of a Stanley-Meyer Cell with Production and Optimization of Brown’s Gas 
      \item Final grade: 1.4
  \end{itemize}

\cvsection{Languages}
  \textbf{Deutsch} \hspace{1.02cm} C1 -- TestDaF \\
  \textbf{Englisch} \hspace{1.05cm} Native \\
  \textbf{Hindi} \hspace{1.45cm} Native

\cvsection{Publication}
  \cvachievement{\faCertificate} 
  {\href{https://www.dgao-proceedings.de/download/125/125_p10.pdf}{Investigation of laser spot size on relative position displacement sensing using laser speckle imaging}}{DGaO - Deutsche Gesellschaft für angewandte Optik, 2024}

\cvsection{Certifications}
  \cvachievement{\faCertificate} 
  {\href{https://www.coursera.org/verify/WYCY2A5EZ6RG}{Introduction to Realtime Embedded Systems Concepts and Practices}}{University of Boulder Colorado, 2026}
  \cvachievement{\faCertificate} 
  {\href{https://www.coursera.org/verify/WYCY2A5EZ6RG}{Linux Systems Programming and Introduction to Buildroot}}{University of Boulder Colorado, 2026}
  \cvachievement{\faCertificate} 
  {\href{https://www.coursera.org/verify/QGONDNKH2CED}{Introduction to Embedded Systems Software and Development Environments}}{University of Boulder Colorado, 2026}
  \cvachievement{\faCertificate} 
  {\href{https://www.coursera.org/verify/specialization/ZUCSRB5PTGA4}{Open Source Software Development}}{The Linux Foundation, 2025}

\cvsection{Awards}
  \cvachievement{\faTrophy}{4 Scholarships}{received in the years 2014-18 for academic performance}
  \divider
  \cvachievement{\faStar}{Medal for the best bachelor’s degree in the 2018 graduating class}{}

% \cvsection{Referenzen}
% References go here.

% \cvsection{Hobbies}
% % Hobbies go here.
% \cvtag{Longboarding}
% \cvtag{Biking}
% \cvtag{Coding} \\
% \cvtag{Cooking}
% \cvtag{Music}
% \cvtag{Hiking}
%
% \cvsection{Interests}
% % Interests go here.
% \cvtag{Robotics}
% \cvtag{IoT}
% \cvtag{Rust}
% \cvtag{Embedded Programming} \\
% \cvtag{Software Automation} \\
% \cvtag{Computer Vision}

\end{paracol}

\end{document}
